% GENERATED FILE. Edit canonical lessons, then run npm run generate:books.
% canonical-source-hash: fnv1a64:6847da00aeae3ff1
% canonical-lessons: KA-C415-amanatu, KA-C415-horagidu, KA-C415-vinayiti, KA-C415-manjuru-madu, KA-C415-adesa

\chapter{Suspensions and Exemptions}
\label{ch:ka-suspensions-and-exemptions}
\hlchaptermodality{415}

\begin{chapteropening}
\textbf{By the end of this chapter:} \emph{I can say suspension, to exclude, an exemption, to sanction and an official order.}

Complete the last lesson of chapter 415: I can say suspension, to exclude, an exemption, to sanction and an official order.
\end{chapteropening}

\section[\texorpdfstring{\kn{ಅಮಾನತು} (amānatu)}{ಅಮಾನತು (amānatu)}]{\kn{ಅಮಾನತು} (amānatu) — suspension}

\label{lesson:KA-C415-amanatu}



\begin{quote}
Before the new one: say the Kannada for to fix, then the Kannada for to restrict.
\end{quote}

\subsection*{You'll want to know: \kn{ಅಮಾನತು}}
\textbf{\kn{ಅಮಾನತು}} — \emph{amānatu} — \textquotedblleft{}suspension\textquotedblright{}.

\begin{quote}
\kn{ಲಂಚ} \kn{ತೆಗೆದುಕೊಂಡ} \kn{ಅಧಿಕಾರಿಯನ್ನು} \kn{ಅಮಾನತು} \kn{ಮಾಡಿದರು}.
\end{quote}

\begin{quote}
\emph{lañca tegedukoṇḍa adhikāriyannu amānatu māḍidaru.}
\end{quote}

\begin{quote}
The official who took a bribe was suspended.
\end{quote}

\kn{ಅಮಾನತು} \kn{ಮಾಡು}, to suspend from a job.

\begin{cousinweb}
\textbf{\kn{ಅಮಾನತು}} (\emph{amānatu}) is usually traced to Arabic \emph{amānat}, \textquotedblleft{}a trust, a deposit\textquotedblright{}, through Hindustani: a post held in trust, not used.
\end{cousinweb}

\begin{grammarlens}[title={What you've built}]
One more word for saying what happens if.
\end{grammarlens}

\subsection*{Your turn}
\begin{itemize}
\raggedright
  \item \emph{Say it:} \emph{amānatu}
  \item \emph{Say it:} \emph{amānatu}, once more
  \item \emph{Say it:} \emph{lañca tegedukoṇḍa adhikāriyannu amānatu māḍidaru.} — The official who took a bribe was suspended.
  \item \emph{Recall:} say the Kannada for to fix, then the Kannada for to restrict, then say \emph{amānatu} again
\end{itemize}

\begin{tcolorbox}[breakable,colback=teal!4,colframe=teal!35!black,title={Before you move on}]
What does \kn{ಅಮಾನತು} mean? (\textquotedblleft{}Suspension\textquotedblright{}.) Say it once more.
\end{tcolorbox}
\section[\texorpdfstring{\kn{ಹೊರಗಿಡು} (horagiḍu)}{ಹೊರಗಿಡು (horagiḍu)}]{\kn{ಹೊರಗಿಡು} (horagiḍu) — to exclude; to leave out}

\label{lesson:KA-C415-horagidu}



\begin{quote}
Before the new one: say the Kannada for to restrict, then the Kannada for suspension.
\end{quote}

\subsection*{You'll want to know: \kn{ಹೊರಗಿಡು}}
\textbf{\kn{ಹೊರಗಿಡು}} — \emph{horagiḍu} — \textquotedblleft{}to exclude; to leave out\textquotedblright{}.

\begin{quote}
\kn{ಗಾಯದಿಂದ} \kn{ಅವನನ್ನು} \kn{ತಂಡದಿಂದ} \kn{ಹೊರಗಿಟ್ಟರು}.
\end{quote}

\begin{quote}
\emph{gāyadinda avanannu taṇḍadinda horagiṭṭaru.}
\end{quote}

\begin{quote}
Because of an injury they left him out of the team.
\end{quote}

Past (I …): \kn{ಹೊರಗಿಟ್ಟೆ}, \emph{horagiṭṭe}.

What is left out takes -\kn{ಅನ್ನು}, the group -\kn{ಇಂದ}.

\begin{cousinweb}
\textbf{\kn{ಹೊರಗಿಡು}} (\emph{horagiḍu}) is native Dravidian, built from \emph{horage}, \textquotedblleft{}outside\textquotedblright{} (DEDR 4333), and \emph{iḍu}, \textquotedblleft{}to put\textquotedblright{} (DEDR 442).
\end{cousinweb}

\begin{grammarlens}[title={What you've built}]
One more word for saying what happens if.
\end{grammarlens}

\subsection*{Your turn}
\begin{itemize}
\raggedright
  \item \emph{Say it:} \emph{horagiḍu}
  \item \emph{Say it:} \emph{horagiḍu}, once more
  \item \emph{Say it:} \emph{gāyadinda avanannu taṇḍadinda horagiṭṭaru.} — Because of an injury they left him out of the team.
  \item \emph{Recall:} say the Kannada for to restrict, then the Kannada for suspension, then say \emph{horagiḍu} again
\end{itemize}

\begin{tcolorbox}[breakable,colback=teal!4,colframe=teal!35!black,title={Before you move on}]
What does \kn{ಹೊರಗಿಡು} mean? (\textquotedblleft{}To exclude; to leave out\textquotedblright{}.) Say it once more.
\end{tcolorbox}
\section[\texorpdfstring{\kn{ವಿನಾಯಿತಿ} (vināyiti)}{ವಿನಾಯಿತಿ (vināyiti)}]{\kn{ವಿನಾಯಿತಿ} (vināyiti) — an exemption; a concession}

\label{lesson:KA-C415-vinayiti}



\begin{quote}
Before the new one: say the Kannada for suspension, then the Kannada for to exclude.
\end{quote}

\subsection*{You'll want to know: \kn{ವಿನಾಯಿತಿ}}
\textbf{\kn{ವಿನಾಯಿತಿ}} — \emph{vināyiti} — \textquotedblleft{}an exemption; a concession\textquotedblright{}.

\begin{quote}
\kn{ಹಿರಿಯರಿಗೆ} \kn{ಶುಲ್ಕದಲ್ಲಿ} \kn{ವಿನಾಯಿತಿ} \kn{ಇದೆ}.
\end{quote}

\begin{quote}
\emph{hiriyarige śulkadalli vināyiti ide.}
\end{quote}

\begin{quote}
Elders get a concession on the fee.
\end{quote}

\kn{ತೆರಿಗೆ} \kn{ವಿನಾಯಿತಿ}, a tax exemption.

\begin{cousinweb}
\textbf{\kn{ವಿನಾಯಿತಿ}} (\emph{vināyiti}) is a Kannada word of the offices whose origin is not settled: neither DEDR nor CDIAL lists it.
\end{cousinweb}

\begin{grammarlens}[title={What you've built}]
One more word for saying what happens if.
\end{grammarlens}

\subsection*{Your turn}
\begin{itemize}
\raggedright
  \item \emph{Say it:} \emph{vināyiti}
  \item \emph{Say it:} \emph{vināyiti}, once more
  \item \emph{Say it:} \emph{hiriyarige śulkadalli vināyiti ide.} — Elders get a concession on the fee.
  \item \emph{Recall:} say the Kannada for suspension, then the Kannada for to exclude, then say \emph{vināyiti} again
\end{itemize}

\begin{tcolorbox}[breakable,colback=teal!4,colframe=teal!35!black,title={Before you move on}]
What does \kn{ವಿನಾಯಿತಿ} mean? (\textquotedblleft{}An exemption; a concession\textquotedblright{}.) Say it once more.
\end{tcolorbox}
\section[\texorpdfstring{\kn{ಮಂಜೂರು} \kn{ಮಾಡು} (mañjūru māḍu)}{ಮಂಜೂರು ಮಾಡು (mañjūru māḍu)}]{\kn{ಮಂಜೂರು} \kn{ಮಾಡು} (mañjūru māḍu) — to sanction; to grant}

\label{lesson:KA-C415-manjuru-madu}



\begin{quote}
Before the new one: say the Kannada for to exclude, then the Kannada for an exemption.
\end{quote}

\subsection*{You'll want to know: \kn{ಮಂಜೂರು} \kn{ಮಾಡು}}
\textbf{\kn{ಮಂಜೂರು} \kn{ಮಾಡು}} — \emph{mañjūru māḍu} — \textquotedblleft{}to sanction; to grant\textquotedblright{}.

\begin{quote}
\kn{ಬ್ಯಾಂಕ್} \kn{ನನಗೆ} \kn{ಸಾಲ} \kn{ಮಂಜೂರು} \kn{ಮಾಡಿತು}.
\end{quote}

\begin{quote}
\emph{byāṅk nanage sāla mañjūru māḍitu.}
\end{quote}

\begin{quote}
The bank granted me a loan.
\end{quote}

Past (I …): \kn{ಮಂಜೂರು} \kn{ಮಾಡಿದೆ}, \emph{mañjūru māḍide}.

\kn{ಮಂಜೂರಾಗು}, to be sanctioned.

\begin{cousinweb}
\textbf{\kn{ಮಂಜೂರು} \kn{ಮಾಡು}} (\emph{mañjūru māḍu}) is built from \emph{mañjūru}, a loan from Hindustani \emph{manzūr}, \textquotedblleft{}approved\textquotedblright{}, which comes from Arabic, with \emph{māḍu}, \textquotedblleft{}to do\textquotedblright{}.
\end{cousinweb}

\begin{grammarlens}[title={What you've built}]
One more word for saying what happens if.
\end{grammarlens}

\subsection*{Your turn}
\begin{itemize}
\raggedright
  \item \emph{Say it:} \emph{mañjūru māḍu}
  \item \emph{Say it:} \emph{mañjūru māḍu}, once more
  \item \emph{Say it:} \emph{byāṅk nanage sāla mañjūru māḍitu.} — The bank granted me a loan.
  \item \emph{Recall:} say the Kannada for to exclude, then the Kannada for an exemption, then say \emph{mañjūru māḍu} again
\end{itemize}

\begin{tcolorbox}[breakable,colback=teal!4,colframe=teal!35!black,title={Before you move on}]
What does \kn{ಮಂಜೂರು} \kn{ಮಾಡು} mean? (\textquotedblleft{}To sanction; to grant\textquotedblright{}.) Say it once more.
\end{tcolorbox}
\section[\texorpdfstring{\kn{ಆದೇಶ} (ādēśa)}{ಆದೇಶ (ādēśa)}]{\kn{ಆದೇಶ} (ādēśa) — an official order}

\label{lesson:KA-C415-adesa}



\begin{quote}
Before the new one: say the Kannada for an exemption, then the Kannada for to sanction.
\end{quote}

\subsection*{You'll want to know: \kn{ಆದೇಶ}}
\textbf{\kn{ಆದೇಶ}} — \emph{ādēśa} — \textquotedblleft{}an official order\textquotedblright{}.

\begin{quote}
\kn{ಸರ್ಕಾರ} \kn{ಶಾಲೆಗಳಿಗೆ} \kn{ಹೊಸ} \kn{ಆದೇಶ} \kn{ಕಳುಹಿಸಿತು}.
\end{quote}

\begin{quote}
\emph{sarkāra śālegaḷige hosa ādēśa kaḷuhisitu.}
\end{quote}

\begin{quote}
The government sent a new order to the schools.
\end{quote}

\kn{ಆದೇಶ} \kn{ಹೊರಡಿಸು}, to issue an order.

\begin{cousinweb}
\textbf{\kn{ಆದೇಶ}} (\emph{ādēśa}) is borrowed from Sanskrit \emph{ādeśa}, \textquotedblleft{}advice, instruction, command\textquotedblright{} (Monier-Williams).
\end{cousinweb}

\begin{grammarlens}[title={What you've built}]
One more word for saying what happens if.
\end{grammarlens}

\subsection*{Your turn}
\begin{itemize}
\raggedright
  \item \emph{Say it:} \emph{ādēśa}
  \item \emph{Say it:} \emph{ādēśa}, once more
  \item \emph{Say it:} \emph{sarkāra śālegaḷige hosa ādēśa kaḷuhisitu.} — The government sent a new order to the schools.
  \item \emph{Recall:} say the Kannada for an exemption, then the Kannada for to sanction, then say \emph{ādēśa} again
\end{itemize}

\begin{tcolorbox}[breakable,colback=teal!4,colframe=teal!35!black,title={Before you move on}]
What does \kn{ಆದೇಶ} mean? (\textquotedblleft{}An official order\textquotedblright{}.) Say it once more.
\end{tcolorbox}
